% GENERATED FILE. Edit canonical lessons, then run npm run generate:books.
% canonical-source-hash: fnv1a64:464e27fcd95660fd
% canonical-lessons: ZH-C479-shangba, ZH-C479-xie, ZH-W479-xie-blood, ZH-C479-zhixue, ZH-C479-xianxue, ZH-C479-xueya

\chapter{Blood}
\label{ch:zh-blood}
\hlchaptermodality{479}

\begin{chapteropening}
\textbf{By the end of this chapter:} \emph{I can talk about a scar, blood, stopping the bleeding, giving blood and blood pressure.}

Complete the last lesson of chapter 479: I can talk about a scar, blood, stopping the bleeding, giving blood and blood pressure.
\end{chapteropening}

\section[shāngbā]{shāngbā — scar}

\label{lesson:ZH-C479-shangba}



\begin{quote}
Before the new one: say the Mandarin for to sprain, then the Mandarin for to dress a wound.
\end{quote}

\subsection*{You'll want to know: shāngbā}
\emph{shāngbā} — \textquotedblleft{}scar\textquotedblright{}.

\begin{quote}
\emph{zhège shāngbā yǐjīng yǒu shí nián le} — I've had this scar for ten years now.
\end{quote}

Its characters come later; for now it is a word for the ear.

\begin{cousinweb}
\emph{shāngbā} is built from \zh{伤} 'to hurt' and \emph{bā} 'a scar'. The first is inherited from Middle Chinese \emph{syang} (Guangyun, in Baxter's transcription), glossed 'to harm'; no dictionary we use records an older form of the second, so its history is not settled.
\end{cousinweb}

\begin{grammarlens}[title={What you've built}]
One more word for saying how it was.
\end{grammarlens}

\subsection*{Your turn}
\begin{itemize}
\raggedright
  \item \emph{Say it:} \emph{shāngbā}
  \item \emph{Say it:} \emph{shāngbā}, once more
  \item \emph{Say it:} \emph{zhège shāngbā yǐjīng yǒu shí nián le} — I've had this scar for ten years now.
  \item \emph{Recall:} say the Mandarin for to sprain, then the Mandarin for to dress a wound, then say \emph{shāngbā} again
\end{itemize}

\begin{tcolorbox}[breakable,colback=teal!4,colframe=teal!35!black,title={Before you move on}]
What does shāngbā mean? (\textquotedblleft{}Scar\textquotedblright{}.) Say it once more.
\end{tcolorbox}
\section[\texorpdfstring{\zh{血} (xiě)}{血 (xiě)}]{\zh{血} (xiě) — blood}

\label{lesson:ZH-C479-xie}



\begin{quote}
Before the new one: say the Mandarin for to dress a wound, then the Mandarin for a scar.
\end{quote}

\subsection*{You'll want to know: \zh{血}}
\textbf{\zh{血}} — \emph{xiě} — \textquotedblleft{}blood\textquotedblright{}.

\begin{quote}
\textbf{\zh{你看这是血}} \emph{nǐ kàn zhè shì xiě} — Look, that's blood.
\end{quote}

The new character, in pieces: \zh{血} — six strokes: a short stroke falling to the left at the top, then a dish of five, as at the foot of \zh{温} \emph{wēn}. The next lesson puts it in your hand.

Alone, and in everyday speech, \zh{血} is said xiě; in words such as liúxuè it is said xuè.

\begin{cousinweb}
\textbf{\zh{血}} (\emph{xiě}) is inherited from Middle Chinese \emph{xwet} (Guangyun, in Baxter's transcription), explained after an older book as what flows out of the flesh. The regular sound change gives \emph{xuè}; the everyday \emph{xiě} is a second, spoken reading, and how it arose is not settled.
\end{cousinweb}

\begin{grammarlens}[title={What you've built}]
One more word for saying how it was.
\end{grammarlens}

\subsection*{Your turn}
\begin{itemize}
\raggedright
  \item \emph{Say it:} \emph{xiě}
  \item \emph{Say it:} \emph{xiě}, once more
  \item \emph{Say it:} \emph{nǐ kàn zhè shì xiě} — Look, that's blood.
  \item \emph{Recall:} say the Mandarin for to dress a wound, then the Mandarin for a scar, then say \emph{xiě} again
\end{itemize}

\begin{tcolorbox}[breakable,colback=teal!4,colframe=teal!35!black,title={Before you move on}]
What does \zh{血} mean? (\textquotedblleft{}Blood\textquotedblright{}.) Say it once more.
\end{tcolorbox}
\section[\texorpdfstring{\zh{血} (xiě)}{血 (xiě)}]{\zh{血} — one short stroke over a dish}

\label{lesson:ZH-W479-xie-blood}



\begin{quote}
The word first, then the character inside it.

\begin{itemize}
\raggedright
  \item \emph{Recall:} say \emph{xiě} — \textbf{R1}, one lesson back, before it decays
\end{itemize}

\begin{quote}
\zh{温} \zh{白}
\end{quote}

\emph{Read it:} each one aloud, with nothing said first — \emph{wēn}, \emph{bái}
\end{quote}

\subsection*{Script — one short stroke over a dish}
\hlblockfigure{\detokenize{figures/ZH-W479-xie-blood-filmstrip.pdf}}{How \zh{血} is written, stroke by stroke}

\begin{quote}
\zh{血}
\end{quote}

\textbf{Two pieces, top to bottom.} Count the pieces, not the strokes.

\textbf{The top, one stroke.} A short stroke falling to the left, above the middle.

\textbf{The dish, five strokes, as at the foot of \zh{温}.} A short vertical on the left, a stroke that runs right and turns down, two short verticals inside, then the long bar across the bottom.

1 + 5 = six.

\subsection*{Your turn}
\begin{itemize}
\raggedright
  \item \emph{Write it:} \zh{血}, all six, with the model in front of you
  \item \emph{Cover:} write it again from memory, and count the pieces, not the strokes
\end{itemize}

\begin{tcolorbox}[breakable,colback=teal!4,colframe=teal!35!black,title={Before you move on}]
How many pieces does \zh{血} have? (\textbf{Two: one stroke and five}.) Where have you written the dish before? (\textbf{At the foot of \zh{温}}.)

Source: \href{https://raw.githubusercontent.com/chanind/hanzi-writer-data/68d10a4b21150cae5e1ebbd223eed289cf32d90c/data/\%E8\%A1\%80.json}{Hanzi Writer Data \zh{血}, pinned stroke paths and medians}.
\end{tcolorbox}
\section[zhǐxuè]{zhǐxuè — to stop the bleeding}

\label{lesson:ZH-C479-zhixue}



\begin{quote}
Before the new one: say the Mandarin for a scar, then the Mandarin for blood.
\end{quote}

\subsection*{You'll want to know: zhǐxuè}
\emph{zhǐxuè} — \textquotedblleft{}to stop the bleeding\textquotedblright{}.

\begin{quote}
\emph{yīshēng yòng le shí fēnzhōng cái zhǐxuè} — It took the doctor ten minutes to stop the bleeding.
\end{quote}

Its characters come later; for now it is a word for the ear.

\begin{cousinweb}
\emph{zhǐxuè} is built from \emph{zhǐ} 'to stop' and \zh{血} 'blood', both inherited from Middle Chinese, \emph{tsyiX} and \emph{xwet} (Guangyun, in Baxter's transcription), which glosses the first 'to stop'.
\end{cousinweb}

\begin{grammarlens}[title={What you've built}]
One more word for saying how it was.
\end{grammarlens}

\subsection*{Your turn}
\begin{itemize}
\raggedright
  \item \emph{Say it:} \emph{zhǐxuè}
  \item \emph{Say it:} \emph{zhǐxuè}, once more
  \item \emph{Say it:} \emph{yīshēng yòng le shí fēnzhōng cái zhǐxuè} — It took the doctor ten minutes to stop the bleeding.
  \item \emph{Recall:} say the Mandarin for a scar, then the Mandarin for blood, then say \emph{zhǐxuè} again
\end{itemize}

\begin{tcolorbox}[breakable,colback=teal!4,colframe=teal!35!black,title={Before you move on}]
What does zhǐxuè mean? (\textquotedblleft{}To stop the bleeding\textquotedblright{}.) Say it once more.
\end{tcolorbox}
\section[xiànxuè]{xiànxuè — to give blood}

\label{lesson:ZH-C479-xianxue}



\begin{quote}
Before the new one: say the Mandarin for blood, then the Mandarin for to stop the bleeding.
\end{quote}

\subsection*{You'll want to know: xiànxuè}
\emph{xiànxuè} — \textquotedblleft{}to give blood\textquotedblright{}.

\begin{quote}
\emph{wǒ měi nián dōu qù xiànxuè} — I give blood every year.
\end{quote}

Its characters come later; for now it is a word for the ear.

\begin{cousinweb}
\emph{xiànxuè} is built from \emph{xiàn} 'to offer' and \zh{血} 'blood', both inherited from Middle Chinese, \emph{xjonH} and \emph{xwet} (Guangyun, in Baxter's transcription). The dictionary lists the first under two readings and glosses \emph{xjonH} 'to present'.
\end{cousinweb}

\begin{grammarlens}[title={What you've built}]
One more word for saying how it was.
\end{grammarlens}

\subsection*{Your turn}
\begin{itemize}
\raggedright
  \item \emph{Say it:} \emph{xiànxuè}
  \item \emph{Say it:} \emph{xiànxuè}, once more
  \item \emph{Say it:} \emph{wǒ měi nián dōu qù xiànxuè} — I give blood every year.
  \item \emph{Recall:} say the Mandarin for blood, then the Mandarin for to stop the bleeding, then say \emph{xiànxuè} again
\end{itemize}

\begin{tcolorbox}[breakable,colback=teal!4,colframe=teal!35!black,title={Before you move on}]
What does xiànxuè mean? (\textquotedblleft{}To give blood\textquotedblright{}.) Say it once more.
\end{tcolorbox}
\section[xuèyā]{xuèyā — blood pressure}

\label{lesson:ZH-C479-xueya}



\begin{quote}
Before the new one: say the Mandarin for to stop the bleeding, then the Mandarin for to give blood.
\end{quote}

\subsection*{You'll want to know: xuèyā}
\emph{xuèyā} — \textquotedblleft{}blood pressure\textquotedblright{}.

\begin{quote}
\emph{wǒ bàba de xuèyā bǐ wǒ gāo} — My dad's blood pressure is higher than mine.
\end{quote}

Its characters come later; for now it is a word for the ear.

\begin{cousinweb}
\emph{xuèyā} is built from \zh{血} 'blood' and \emph{yā} 'to press', both inherited from Middle Chinese, \emph{xwet} and \emph{'aep} (Guangyun, in Baxter's transcription), which glosses the second 'to press down'.
\end{cousinweb}

\begin{grammarlens}[title={What you've built}]
One more word for saying how it was.
\end{grammarlens}

\subsection*{Your turn}
\begin{itemize}
\raggedright
  \item \emph{Say it:} \emph{xuèyā}
  \item \emph{Say it:} \emph{xuèyā}, once more
  \item \emph{Say it:} \emph{wǒ bàba de xuèyā bǐ wǒ gāo} — My dad's blood pressure is higher than mine.
  \item \emph{Recall:} say the Mandarin for to stop the bleeding, then the Mandarin for to give blood, then say \emph{xuèyā} again
\end{itemize}

\begin{tcolorbox}[breakable,colback=teal!4,colframe=teal!35!black,title={Before you move on}]
What does xuèyā mean? (\textquotedblleft{}Blood pressure\textquotedblright{}.) Say it once more.
\end{tcolorbox}
